% from ega4-2-fr.tex
\oldpage[IV-2]{23}
notons que la question est locale sur $X$, et l'on peut donc supposer $X$
affine~; on sait alors (2.2.12) qu'il y a un schéma affine $X''$ et un
morphisme fidèlement plat $g:X''\to X'$ qui soit un isomorphisme local~; par
suite, il revient au même de dire que $f^*(\mathscr F)$ possède la propriété
$\mathbf P$, ou que $g^*(f^*(\mathscr F))$ la possède. On est ainsi ramené au
cas où $X=\operatorname{Spec}(A)$, $X'=\operatorname{Spec}(A')$~; compte tenu
de (2.2.3) et (II, 6.1.4.1), il suffit donc de prouver le

\begin{lemma}[2.5.3]
\label{IV.2.5.3-fr}
Soient $A$ un anneau, $A'$ une $A$-algèbre fidèlement plate, $M$ un
$A$-module, $M'=M\otimes_A A'$. Pour que $M$ soit de type fini (resp. de
présentation finie), il faut et il suffit que $M'$ le soit.
\end{lemma}

Pour la démonstration, voir Bourbaki, \emph{Alg. comm.},
chap.~I\textsuperscript{er}, \S~3, n\textsuperscript{o}~6, prop.~11.

\begin{remark}[2.5.4]
\label{IV.2.5.4-fr}
Les assertions de (2.5.2) pour les propriétés (i) et (ii) sont encore
valables si l'on suppose seulement $f$ quasi-fidèlement plat (2.3.3) et
quasi-compact. On est en effet ramené (Bourbaki, \emph{loc. cit.}) à prouver
le
\end{remark}

\begin{lemma}[2.5.4.1]
\label{IV.2.5.4.1-fr}
Soit $\rho:A\to A'$ un homomorphisme d'anneaux, tel que le morphisme
correspondant $f:\operatorname{Spec}(A')\to\operatorname{Spec}(A)$ soit
surjectif~; supposons qu'il existe un $A'$-module $N'$ de type fini, qui soit
$A$-plat et ait pour support $\operatorname{Spec}(A')$. Alors, si
$u:P\to Q$ est un homomorphisme de $A$-modules tel que
\[
u\otimes 1:P\otimes_A A'\longrightarrow Q\otimes_A A'
\]
soit surjectif, $u$ est surjectif.
\end{lemma}

En effet, on en déduit d'abord que l'homomorphisme
\[
u\otimes 1_{N'}:(P\otimes_A A')\otimes_{A'}N'
\longrightarrow(Q\otimes_A A')\otimes_{A'}N'
\]
est surjectif. Soit $\mathfrak q$ un idéal premier de $A'$, et soit
$\mathfrak p=\rho^{-1}(\mathfrak q)$~; l'homomorphisme
$(P\otimes_A N')_{\mathfrak q}\to(Q\otimes_A N')_{\mathfrak q}$
correspondant est surjectif, et il s'écrit
\[
u_{\mathfrak p}\otimes 1:
P_{\mathfrak p}\otimes_{A_{\mathfrak p}}N'_{\mathfrak q}
\longrightarrow
Q_{\mathfrak p}\otimes_{A_{\mathfrak p}}N'_{\mathfrak q}
\quad(0_{\mathrm I},1.5.4).
\]
Par hypothèse on a $N'_{\mathfrak q}\ne0$ et $N'_{\mathfrak q}$ est un
$A_{\mathfrak p}$-module plat $(0_{\mathrm I},6.3.1)$~; en vertu du lemme de
Nakayama, on a
$\mathfrak qN'_{\mathfrak q}\ne N'_{\mathfrak q}$, et \emph{a fortiori}
$\mathfrak pN'_{\mathfrak q}\ne N'_{\mathfrak q}$, donc
$N'_{\mathfrak q}$ est un $A_{\mathfrak p}$-module fidèlement plat
$(0_{\mathrm I},6.4.1)$. Il en résulte que $u_{\mathfrak p}$ est surjectif
$(0_{\mathrm I},6.4.1)$, et ceci ayant lieu pour tout
$\mathfrak p\in\operatorname{Spec}(A)$, puisque $f$ est surjectif, on en
conclut finalement que $u$ est surjectif (Bourbaki, \emph{Alg. comm.},
chap.~II, \S~3, n\textsuperscript{o}~3, th.~1).

\begin{proposition}[2.5.5]
\label{IV.2.5.5-fr}
Soient $f:X\to Y$ un morphisme, $\mathscr F$ un $\mathcal O_X$-Module
quasi-cohérent de type fini et $f$-plat, $\mathscr E$ un
$\mathcal O_Y$-Module quasi-cohérent de type fini~; pour tout $y\in Y$, on
pose $\mathscr F_y=\mathscr F\otimes_{\mathcal O_Y}k(y)$. Pour qu'un point
$x\in X$ soit point maximal de
$\operatorname{Supp}(\mathscr E\otimes_Y\mathscr F)$, il faut et il suffit que
$y=f(x)$ soit point maximal de $\operatorname{Supp}(\mathscr E)$, et que $x$
soit point maximal de $\operatorname{Supp}(\mathscr F_y)$ dans $f^{-1}(y)$.
S'il en est ainsi, on a
\end{proposition}

\begin{env}[2.5.5.1]
\label{IV.2.5.5.1-fr}
\[
\operatorname{long}\bigl((\mathscr E\otimes_Y\mathscr F)_x\bigr)
=\operatorname{long}(\mathscr E_y)\,
 \operatorname{long}\bigl((\mathscr F_y)_x\bigr).
\]
\end{env}

Il est clair que
$f(\operatorname{Supp}(\mathscr E\otimes_Y\mathscr F))
\subset\operatorname{Supp}(\mathscr E)$ $(0_{\mathrm I},5.2.2)$~; l'image
par $f$ de toute composante irréductible de
$\operatorname{Supp}(\mathscr E\otimes_Y\mathscr F)$ est donc contenue dans
une composante irréductible de $\operatorname{Supp}(\mathscr E)$. On peut se
borner au cas où $\operatorname{Supp}(\mathscr F)=X$. En effet
(\textbf{Err}\textsubscript{\textbf{III}}, 30), puisque $\mathscr F$ est de
type fini, il existe un sous-préschéma fermé $X'$ de $X$ ayant
$\operatorname{Supp}(\mathscr F)$ pour espace sous-jacent et un
$\mathcal O_{X'}$-Module quasi-cohérent de type fini $\mathscr G$ tels que,
si $j:X'\to X$ est l'injection canonique, on ait
$\mathscr F=j_*(\mathscr G)$. Si l'on pose $f'=f\circ j$, il est clair que
$\mathscr G$ est $f'$-plat et que l'on a
\[
\operatorname{Supp}(\mathscr E\otimes_Y\mathscr F)
=\operatorname{Supp}(\mathscr E\otimes_Y\mathscr G).
\]
Supposons donc $\operatorname{Supp}(\mathscr F)=X$~; si $Z$ est une
composante irréductible de $\operatorname{Supp}(\mathscr E)$, on a
\[
f^{-1}(Z)\subset\operatorname{Supp}(\mathscr E\otimes_Y\mathscr F)
\quad\text{(I, 9.1.13),}
\]
et il résulte de (2.3.4) que toute composante irréductible de $f^{-1}(Z)$
domine $Z$~; autrement dit, si $x$ est point générique d'une composante
irréductible de $\operatorname{Supp}(\mathscr E\otimes_Y\mathscr F)$
contenue dans $f^{-1}(Z)$, $y=f(x)$ est le point générique de $Z$~; en outre
$(0_{\mathrm I},2.1.8)$, $x$ est point générique d'une composante
irréductible de
$f^{-1}(y)=\operatorname{Supp}(\mathscr F_y)$ (I, 9.1.13), et
réciproquement tout point générique d'une de ces composantes est point
générique d'une composante irréductible de $f^{-1}(Z)$.

